\originalsection{18.315：作业7（原卷 Fall 2005）}
本组收录原题3的三小问. 原题1、2的排列规避与队列/栈,
原题4的多边形台球轨道另登记为本册范围外内容.

\begin{exercise}\label{mit:18315-hw7-q3}
\textbf{18.315, 原卷 Fall 2005, 作业7, 题3.}
有限简单图 $H$（至少两个顶点）的密度为
$\alpha(H)=|E(H)|/\binom{|V(H)|}{2}$.
对无限简单图 $G$, 令 $\alpha_n(G)$ 为全部 $n$ 顶点诱导子图的最大密度,
$\alpha(G)=\limsup_{n\to\infty}\alpha_n(G)$.
证明: (a) 每个 $a\in[0,1]\cap\mathbb Q$ 是某个有限图的密度;
(b) 无限图的 $\alpha(G)\ne1/3$; (c) 无限图的 $\alpha(G)$ 总是有理数.
\end{exercise}
\sourceof{官方 HW7 第1页题3; 原页(b)是“不等于 $1/3$”, 文本提取丢失了不等号; 编者独立解答}
\begin{solution}
(a) 写 $a=p/q$, $0\le p\le q$, $q\ge1$.
取 $2q$ 个顶点, 候选边有 $q(2q-1)$ 条,
任选其中 $p(2q-1)$ 条, 密度恰为 $p/q$.

先说明无限图的极限实际存在. 在一个实现 $\alpha_{n+1}(G)$ 的诱导子图中随机删除一个顶点,
剩余子图的平均密度等于删除前的密度, 因每条边以概率 $(n-1)/(n+1)$ 留下.
至少一个剩余子图的密度不小于平均值, 故 $\alpha_n(G)\ge\alpha_{n+1}(G)$.
最大值确实能取得, 因边数只有 $0,1,\ldots,\binom n2$ 这些有限种整数取值.
所以 $\alpha_n(G)$ 非增且有下界0, 收敛到 $\alpha(G)$.

下面证明所需的有限图吹胀引理:
对任意固定 $r,t\ge1$ 与 $\delta>0$,
足够大的有限图若密度至少 $1-1/r+\delta$,
就含有每部 $t$ 个顶点的完全 $(r+1)$ 部图.
记 $p=1-1/r$, 只需考虑 $p+\delta\le1$.
选固定整数 $h\ge r+1$ 足够大, 使
$p h^2/2\le(p+\delta/2)\binom h2$.
均匀抽取 $h$ 个顶点, 平均边数至少 $(p+\delta)\binom h2$.
若抽出的图无 $K_{r+1}$, Turán 定理把它的边数限制在 $ph^2/2$ 以下;
其他图的边数最多 $\binom h2$.
因此至少 $\delta/2$ 比例的 $h$ 顶点集含一个 $K_{r+1}$.
一份 $K_{r+1}$ 至多被 $\binom{N-r-1}{h-r-1}$ 个 $h$ 顶点集计入,
于是原 $N$ 顶点图中的 $(r+1)$ 团数至少
\[
\frac{\delta}{2}\,
\frac{\binom N{r+1}}{\binom h{r+1}}
\ge cN^{r+1}
\]
（$N$ 足够大, $c>0$ 固定）.

还需证明稠密超图含完整多部结构, 才能把这些团吹胀.
我们证: 对固定 $k,t,c>0$, 一个每部大小至多 $N$、有至少 $cN^k$ 条横跨各部的边的
$k$ 部 $k$ 一致超图, 在 $N$ 足够大时含 $K^{(k)}_{t,\ldots,t}$.
对 $k$ 归纳. $k=1$ 是至少 $cN$ 个点中选 $t$ 个.
设 $k\ge2$, 对其余 $k-1$ 部的每个顶点组 $e$, 记第一部中的共同邻点数为 $d_e$.
有 $\sum_e d_e\ge cN^k$, 每项至多 $N$, 项数至多 $N^{k-1}$;
所以至少 $(c/2)N^{k-1}$ 个 $e$ 满足 $d_e\ge(c/2)N$.
计数这些 $e$ 与它们的 $t$ 元共同邻点集 $S$ 的配对, 当 $N\ge4t/c$ 时,
\[
\sum_e\binom{d_e}{t}
\ge\frac c2 N^{k-1}\frac{(cN/4)^t}{t!}.
\]
第一部中 $S$ 至多有 $N^t/t!$ 种,
故某个 $S$ 的共同链超图有至少
$(c/2)(c/4)^tN^{k-1}$ 条边.
应用归纳假设, 在其余各部各选 $t$ 个点, 与 $S$ 一起得到完整超图.
若原超图不是多部的, 随机把顶点分到 $k$ 部,
每条 $k$ 元边以固定概率 $k!/k^k$ 横跨各部,
所以某个划分仍保留固定正比例的边, 同样可用此引理.

将有限图的 $(r+1)$ 团作为超边, 上述引理给出
$K^{(r+1)}_{t,\ldots,t}$.
取其中任意两个不同部分的点, 可从其他部分各添一点成为超边,
所以这两点在原图中相邻. 因而原图含所需完整 $(r+1)$ 部图,
有限图吹胀引理证明完成.

(b),(c) 现在可一起得到更精确的分类.
设 $\alpha=\alpha(G)<1$, 选唯一 $r\ge1$ 使
\[
1-\frac1r\le\alpha<1-\frac1{r+1}.
\]
若左侧为严格不等式, 取
$0<\delta<\alpha-(1-1/r)$.
对任意 $t$, 足够大的 $n$ 顶点最大密度诱导子图均满足吹胀引理,
故 $G$ 含每部 $t$ 点的完整 $(r+1)$ 部图.
它所在的诱导子图还可有同部边, 只会使密度增加,
因此对 $v=(r+1)t$ 有
\[
\alpha_v(G)\ge\left(1-\frac1{r+1}\right)\frac v{v-1}.
\]
令 $t\to\infty$ 得 $\alpha\ge1-1/(r+1)$, 矛盾.
所以必有 $\alpha=1-1/r$. 加上 $\alpha=1$ 的情形,
全部可能值只能属于
$\{0,1/2,2/3,3/4,\ldots\}\cup\{1\}$.
它们都是有理数, 且没有 $1/3$, 即证明(b),(c).
每个列出的值也确能实现: 用各部无限的完整 $r$ 部图, 以及无限完全图;
在有限诱导子图中均衡选各部顶点即可逼近相应密度.
\end{solution}
